\section{Modelo de privacidad y permisos: cómo debe restringirse al doble digital}

Para que un doble digital salga al mundo de forma proactiva en nombre del usuario, necesita un modelo de permisos. Qué context puede consultar un doble, con quién puede hablar, qué puede expresar en representación del usuario, qué acciones puede completar de forma automática y cuáles requieren confirmación: todas son cuestiones de fondo del producto. Sin un modelo de permisos, cuanto más proactivo sea el doble, más peligroso resulta.

\subsection{Cuatro niveles de permisos}

\noindent\begin{tabularx}{\textwidth}{p{0.22\textwidth}p{0.34\textwidth}X}
\toprule
Nivel de permisos & Qué puede hacer el doble & Control de riesgos \\
\midrule
Nivel de lectura & Leer los datos del usuario, sus expresiones previas, sus preferencias y sus relaciones & El usuario puede consultar, eliminar y autorizar por niveles. \\
Nivel de borrador & Generar comentarios, mensajes privados y justificaciones de recomendaciones, pero sin enviarlos & Se publica tras la confirmación del usuario. \\
Nivel de acciones de bajo riesgo & Dar «me gusta», guardar en favoritos, interacciones públicas de bajo riesgo & Límites de frecuencia y registros de auditoría. \\
Nivel de acciones de alto riesgo & Conversaciones privadas, opiniones sensibles, compromisos en las relaciones, acciones comerciales & Requiere confirmación humana obligatoria y debe poder revertirse. \\
\bottomrule
\end{tabularx}

\begin{knowledgebox}{El modelo de permisos también es experiencia de producto}
La condición para que el usuario esté dispuesto a entregar su context es saber que el doble no se saldrá de control. Unos permisos claros no limitan el producto: hacen que el usuario se atreva a usarlo.
\end{knowledgebox}

\subsection{La propiedad del context}

La cuestión más delicada de las redes sociales con IA es a quién pertenece realmente el context. Las conversaciones, las obras, las preferencias sociales y los datos de relaciones del usuario no pueden ser solo un activo de la plataforma. A largo plazo, el usuario necesita derechos sobre su context para poder trasladarlo, exportarlo, eliminarlo y guardarlo localmente. De lo contrario, cuanto mejor conozca la plataforma al usuario, más inquieto se sentirá este.

\begin{importantbox}{Lista de derechos sobre el context}
El usuario debería tener, como mínimo, derecho de consulta, derecho a una explicación, derecho de eliminación, derecho de exportación, derecho de autorización por niveles y derecho a confirmar las acciones sensibles. Si una red social con IA quiere construir confianza a largo plazo, debe convertir estos derechos en capacidades del producto.
\end{importantbox}

\subsection{Resumen de la sección}

Cuanto más proactivo es el doble digital, más importante se vuelve el modelo de permisos. La base de la confianza en las redes sociales con IA no es que el modelo hable mejor, sino que el usuario pueda controlar su propio context y el comportamiento de su doble.
